Un framework es un conjunto de herramientas que facilita el desarrollo de 
aplicaciones. Proporciona una base sólida y reutilizable para crear aplicaciones 
más rápidas, robustas y eficientes.

Para el desarrollo de este proyecto se han valorado diferentes frameworks con el 
fin de buscar un desarrollo eficiente y seguro. Antes de analizar las opciones disponibles, 
es importe destacar que una aplicación web consta de dos partes fundamentales \cite{laboratoria2021}: 

\begin{description}
    \item [Backend] Se encarga de la lógica, la gestión de datos y la conexión a la base de datos. Esta parte no es visible para el usuario final
    \item [Frontend] Es la interfaz con la que interactúa el usuario. Incluye elementos como texto, botones e imágenes con el fin de facilitar su navegación. 
\end{description}
